\begin{resumo}[Abstract]
  \begin{otherlanguage*}{english}
    Health service scheduling in Brazil occurs predominantly
    by phone or in person: research by Doctoralia and TuoTempo
    (2022) indicates that 91\% of medical centers use the phone
    as their primary scheduling channel, with one in every three
    calls going unanswered. This inefficiency creates difficulties
    for both patients, who depend on business hours to book
    appointments, and health professionals and clinics, which
    manage their schedules through spreadsheets and messaging
    applications. In this context, this work proposes the
    development of Medicano, a web-based health service
    scheduling platform that connects patients to clinics and
    professionals across different specialties, including
    medicine, psychology, psychiatry, dentistry, and nutrition.
    The platform incorporates an intelligent triage chatbot
    based on Large Language Models that, through a guided
    conversation, identifies the user's needs and directs them
    to the most suitable specialty. The solution was developed
    as a monorepo using React and TypeScript on the frontend,
    Node.js with NestJS on the backend, MongoDB as the NoSQL
    database, and Vercel AI SDK for language model integration.
    The infrastructure is composed of a frontend hosted on
    Amazon S3 with CloudFront distribution and a backend
    running on an Amazon EC2 instance with Nginx as a reverse
    proxy, with secrets managed by AWS Secrets Manager.
    Development followed an agile methodology with two-week
    sprints, resulting in a minimum viable product covering
    the flows of registration and authentication, intelligent
    triage by chatbot, search for professionals and clinics
    by specialty and location, appointment scheduling, and
    agenda management. This work demonstrates the technical
    feasibility of the platform and establishes a solid
    foundation for its continuous evolution.

    \vspace{\onelineskip}
    \noindent\textbf{Keywords:}
    health services scheduling. intelligent triage.
    large language models. health marketplace.
    cloud computing.
  \end{otherlanguage*}
\end{resumo}